% 引言
\section{引言}

数字经济已成为全球经济结构转型的核心驱动力。截至 2023 年，中国数字经济规模超过 50 万亿元，占 GDP 比重超过 40\%，连续多年位居世界第二。与此同时，中国就业市场正经历深刻的结构性变革：第三产业就业占比从 2011 年的 35.7\% 持续上升至 2023 年的 49.1\%，超越第二产业成为吸纳就业的主体；传统制造业岗位加速收缩，平台经济、数字服务、零工经济等新业态不断创造新的就业岗位，就业形态日益多元化。在此宏观背景下，一个核心政策问题浮现：数字经济发展究竟如何影响省际就业结构？其效应在不同城镇化水平和区域之间是否呈现非线性特征？

在理论上，数字经济对就业存在两种方向相反的力量。Acemoglu 和 Restrepo（2018, 2019）提出的任务框架（Task Framework）将自动化建模为资本在生产任务中替代劳动的过程，同时引入新任务创造作为劳动需求的补偿机制。Autor（2015）从历史视角论证，技术变革一贯改变工作内容而非消灭工作总量——ATM 机没有消灭银行柜员，电子表格软件没有消灭会计职业。这两种力量——替代效应与补偿效应——的净结果在理论上并不确定，取决于技能结构、制度环境和观察的时间窗口。这一理论上的不确定性为中国情境的实证研究提供了天然的学术张力。

已有关于中国数字经济就业效应的实证研究主要存在以下不足。第一，多数文献的数据窗口截至 2020 年（赵涛等，2020；王芳，2023；赵向豪等，2024），未能覆盖 COVID-19 后数字经济加速期（2020--2023 年）对就业结构的影响。疫情三年间，数字经济渗透率发生了质的跃升——在线办公、远程教育、数字医疗等场景的普及速度远超此前任何时期。这一"自然加速"如何影响就业结构，现有文献缺乏系统检验。第二，现有研究以线性模型为主。少数研究（张颖和辛爱洪，2024；李明和张强，2023）开始探讨门槛效应，但对城镇化水平如何调节数字经济就业效应的非线性特征探讨有限。城镇化作为数字基础设施落地的物理条件和经济密度条件，理论上应在数字经济就业效应中发挥调节作用，但这一假说尚未得到充分检验。第三，多数研究使用单一数字金融指数作为数字经济的代理变量。对广义数字经济综合指数（覆盖互联网基础设施、移动通信、软件服务业等维度）与狭义数字金融指数的比较分析不足，导致无法判断"数字化的哪个维度"在推动就业结构转型中发挥主导作用。

本文试图在上述三个方面做出边际贡献。第一，将数据窗口延长至 2023 年，覆盖数字经济加速期，估计其对省际就业结构的最新效应。第二，采用 Hansen（1999）面板门槛模型，以城镇化率为门槛变量，检验数字经济就业结构效应的非线性特征。第三，同时使用自建数字经济综合指数和北大数字普惠金融指数（郭峰等，2020），通过双指标对比揭示广义数字化转型与狭义金融科技在就业效应上的差异，为政策优先序提供经验依据。

本文其余部分安排如下：第二部分为文献综述与理论假说，梳理技术变革与就业的理论脉络、数字经济测度方法、已有实证证据，并在此基础上提出三个研究假说；第三部分介绍模型设定、变量构造与数据来源；第四部分报告基准回归、门槛效应与区域异质性分析结果；第五部分从替换核心解释变量、剔除直辖市、对数变换、剔除 COVID 年份和替代缩尾水平五个维度进行稳健性检验；第六部分总结发现并讨论政策含义与研究局限。
